% !TEX program = pdflatex
% !TEX encoding = UTF-8 Unicode
\documentclass[11pt,a4paper,oneside]{report}

\usepackage[english]{babel}
\usepackage[utf8]{inputenc}
\usepackage[T1]{fontenc}
\usepackage{lmodern}
\usepackage{amsmath,amssymb,amsthm}
\usepackage{geometry}
\usepackage{graphicx}
\usepackage{booktabs}
\usepackage{longtable}
\usepackage{siunitx}
\usepackage{hyperref}
\usepackage{listings}
\usepackage{xcolor}
\usepackage{fancyhdr}
\usepackage{setspace}

\geometry{margin=2.5cm}
\graphicspath{{figures/}}

\lstset{
  language=Bash,
  basicstyle=\ttfamily\small,
  frame=single,
  numbers=left,
  numbersep=5pt,
  backgroundcolor=\color{yellow!10},
  breaklines=true,
  columns=fullflexible,
}

\title{Documentation of the HIF--Iron Homeostasis Simulation Code}
\author{Irina Piazza\\
  with the original framework by Eric Fanchon and Alexandre Rocca\\[1em]
  \normalsize TIMC-IMAG, CNRS, Grenoble, France}
\date{}

\begin{document}
\maketitle

\begin{abstract}
This report documents the Python simulation framework used to explore the
parameter space of a cellular iron homeostasis model and its extension toward
HIF--iron coupling. It describes the software architecture of the public
release, how to install and run the code, the sampling and filtering workflow,
and selected technical notes (including stiffness diagnostics). Historical
laboratory protocols are retained in corrected English for reproducibility of
earlier campaigns; command names are updated to the current public API
(\texttt{simulate.py}, \texttt{config.py}, \texttt{filter.py}, \texttt{model.py}).
\end{abstract}

\tableofcontents
\newpage

\chapter{Overview}

\section{Scientific context}
Iron homeostasis is tightly regulated at the cellular level. The models studied
in this project describe key species (labile iron, ferritin and iron storage,
IRP1, TfR1, ferroportin, and later HIF) and their kinetic interactions under
different iron-replenishment conditions. The numerical goal is to sample a
high-dimensional parameter space, retain only trajectories that satisfy
biological constraints, and characterise the resulting admissible parameter
regions.

\section{Software lineage}
The original simulation framework was developed by Eric Fanchon and Alexandre
Rocca (TIMC-IMAG, CNRS). Irina Piazza later reorganised the code for clarity and
maintainability: configuration was centralised, input/output was standardised
with pandas DataFrames, filtering and logging were rewritten, plotting utilities
were added, and the ODE right-hand side was unified. The public release
distributed with this documentation is the HIF--iron oriented version of that
work, modernised for Python~3.

Changes introduced by Irina Piazza are marked in the source with comments of the
form \texttt{\# IP <date>}.

\section{Architecture of the public code}
The main modules are:

\begin{description}
  \item[\texttt{simulate.py}] Entry point. Implements sampling modes
  (\texttt{Box-Sampling}, \texttt{Re-Sampling}, \texttt{Shell-Sampling},
  \texttt{PCA-guided-sampling}, \ldots), numerical integration, constraint
  checks, optional filtering, dumping of trajectories, and plotting.
  \item[\texttt{model.py}] Model-specific ODE system, algebraic steady-state
  construction, absolute/relative constraints, and optional stiffness
  diagnostics.
  \item[\texttt{config.py}] Global constants: solver tolerances, sampling
  settings, variable/parameter names, bounds and biological constraints.
  \item[\texttt{filter.py}] Post-processing filter on ratios and concentration
  bounds; writes filtered solution lists and enclosing boxes.
  \item[\texttt{read\_file.py}] Reads bound files and results tables.
  \item[\texttt{log.py}] Writes \texttt{info.txt} and copies input files into
  each run directory.
  \item[\texttt{sobol\_lib.py}, \texttt{pca.py}] Quasi-Monte Carlo sampling and
  PCA-guided resampling helpers.
  \item[\texttt{plot\_*.py}, \texttt{multiplot\_*.py}] Offline visualisation
  utilities.
\end{description}

Each execution creates a timestamped folder \texttt{Run\_DD\_MM\_YYYY-HhMmSs/}
containing results, plots, filtered outputs (when applicable), and run metadata.

\section{Typical pipeline}
\begin{enumerate}
  \item Sample parameter vectors inside a box (Sobol sequence) or around
  previous solutions (re-sampling / shell / PCA modes).
  \item For each vector, compute algebraic initial conditions, improve the
  steady state numerically if needed, then integrate the ODEs for each iron
  input listed in \texttt{MULTIPLE\_INPUTS}.
  \item Accept the set only if absolute and relative constraints are satisfied.
  \item After the sampling loop, filter the accepted set with
  \texttt{filter.py}, dump a subset of trajectories, and produce plots.
\end{enumerate}

\chapter{Installation and usage}

\section{Environment setup}
Create an isolated Python environment and install dependencies from the project
root:

\begin{lstlisting}
python -m venv .venv
# Linux/macOS:
source .venv/bin/activate
# Windows PowerShell:
# .\.venv\Scripts\Activate.ps1

python -m pip install --upgrade pip
pip install -r requirements.txt
python tests/run_all_tests.py
\end{lstlisting}

Required packages include NumPy, SciPy, pandas, matplotlib and recordclass
(see \texttt{requirements.txt}).

\section{Running a simulation}
Basic command:

\begin{lstlisting}
python simulate.py b-vars-hif.txt b-par-hif.txt results.txt Box-Sampling 1000
\end{lstlisting}

Arguments:
\begin{enumerate}
  \item Variable bounds file (\texttt{name lower\_bounds upper\_bounds}).
  \item Parameter bounds file, or a previous results / filter-solution file
  depending on the mode.
  \item Output results file name (written under the new \texttt{Run\_*} folder).
  \item Mode (optional; default is box sampling when the CLI provides five
  arguments as above).
  \item Number of sampled points / iterations.
\end{enumerate}

Useful modes:
\begin{description}
  \item[\texttt{Box-Sampling}] Sample the parameter box with a Sobol sequence.
  \item[\texttt{Re-Sampling}] Build local boxes around previously accepted
  solutions and sample them.
  \item[\texttt{Shell-Sampling}] Sample a shell of controlled width around a
  previous solution set (\texttt{magnitude} in \texttt{config.py}).
  \item[\texttt{PCA-guided-sampling}] Resample using a PCA reduction of the
  accepted cloud.
\end{description}

Key settings live in \texttt{config.py}: solver tolerances
(\texttt{SIMU\_RTOL}, \texttt{SIMU\_ATOL}), steady-state criteria, iron inputs
(\texttt{MULTIPLE\_INPUTS}), biological bounds, and filter ratios.

\section{Outputs}
A typical run directory contains:
\begin{itemize}
  \item the results table (parameters, initial variables, ratios, total iron);
  \item \texttt{Solutions\_y\_t/} and \texttt{Solutions\_filtered\_y\_t/}
  trajectory dumps (a capped subset, to keep files manageable);
  \item plot folders for raw and filtered solutions, and parameter correlations;
  \item \texttt{Filter\_sol\_*.txt} / \texttt{Filter\_box\_*.txt} when filtering
  succeeds;
  \item \texttt{info.txt} with settings and copied input bound files.
\end{itemize}

\section{Regression tests}
Contract and numerical fingerprint tests protect against unintentional changes:

\begin{lstlisting}
python tests/run_all_tests.py
\end{lstlisting}

Golden cases are stored under \texttt{tests/golden/}. Regenerate them only after
a deliberate, reviewed numerical change:

\begin{lstlisting}
python tests/generate_golden.py
\end{lstlisting}

\chapter{Historical sampling protocol}
\label{chap:protocol}

This chapter summarises laboratory campaigns performed while refining the
iron-homeostasis parameter region (late 2019 -- early 2020). The narrative is
preserved for reproducibility of those results. In the present public code, the
entry point is \texttt{simulate.py} (formerly \texttt{simulate\_test.py} /
\texttt{simulate\_P.py}), configuration lives in \texttt{config.py}, and
filtering is integrated after sampling (formerly a separate
\texttt{filter\_test.py} call in early workflows).

\section{Integrated simulate + filter workflow}
In later versions of the code, filtering is invoked in the same run as sampling:
once \texttt{apply\_method()} returns candidate solutions, \texttt{filter()}
retains only those that satisfy the ratio and concentration constraints. The
code also writes trajectory files and plots for each variable and each iron
experiment. With seven variables and four iron inputs this yields up to
28 trajectory files and 28 plots for both the unfiltered and filtered sets.

Because the number of accepted solutions can be large, trajectory dumping and
plotting are capped (on the order of 100 trajectories written and a smaller
number plotted) so that output remains manageable. The main results table stores
parameters, initial variables, ratios (Fe-sto, TfR1, Ft) and total iron for each
experiment in a DataFrame-like columnar layout.

\section{Outline of the refinement campaign}
A representative sequence was:
\begin{enumerate}
  \item \textbf{Box sampling} of a large initial box (millions of Sobol points)
  with relatively relaxed ferritin-ratio bounds at 250\,nM iron replenishment
  (e.g.\ Ft ratio in $[4,12]$).
  \item \textbf{Filtering}, which may initially retain very few or no points.
  \item Repeated \textbf{re-sampling} around the surviving solutions, alternating
  with filtering, to densify the admissible region.
  \item Tightening of the ferritin ratio at 250\,nM to a validated interval
  such as $[4, 6.5]$.
  \item Optional \textbf{shell sampling} to explore the boundary of the current
  box, followed again by re-sampling / box sampling to restore a more uniform
  coverage.
  \item Comparison of the enclosing box volume with the reference box obtained
  by A.~Rocca, and occasional enlargement of weakly constrained directions
  (e.g.\ $k_{\mathrm{IRP\_FPN1a}}$) when biological arguments suggested a wider
  admissible range.
\end{enumerate}

Table~\ref{tab:protocol-summary} gives a compact schematic of one such campaign
(counts are those recorded at the time; exact folder names varied between runs).

\begin{table}[!ht]
\centering
\begin{tabular}{clccc}
\toprule
Step & Mode & Iterations & Ft ratio at 250\,nM & Filtered sols. \\
\midrule
1 & Box-Sampling & $4\times10^6$ & $[4.0,12.0]$ & 0 after filter \\
2A & Re-Sampling ($0.2$) & $10^3$ & $[4.0,12.0]$ & 1 \\
2B & Re-Sampling ($0.2$) & $10^5$ & $[4.0,12.0]$ & $12421$ \\
2C & Re-Sampling & (dense) & $[4.0,6.5]$ & $\sim 922$ \\
later & Shell / Box / Re-Sampling & up to $5\times10^6$ & $[4.0,6.5]$ & variable \\
\bottomrule
\end{tabular}
\caption{Schematic of the historical refinement steps. Intermediate unfiltered
counts were typically much larger than the filtered counts.}
\label{tab:protocol-summary}
\end{table}

\section{Example commands (updated names)}
Initial box sampling:

\begin{lstlisting}
python simulate.py b-vars-hif.txt b-par-hif.txt results.txt Box-Sampling 4000000
\end{lstlisting}

Re-sampling around previous solutions (use a results or filter-solution file as
the second argument):

\begin{lstlisting}
python simulate.py b-vars-hif.txt previous_Filter_sol.txt results_rs.txt Re-Sampling 100000
\end{lstlisting}

Shell sampling:

\begin{lstlisting}
python simulate.py b-vars-hif.txt previous_Filter_sol.txt results_shell.txt Shell-Sampling 4000000
\end{lstlisting}

\section{Comparison with A.~Rocca's box}
At intermediate stages the enclosing volume of the newly obtained box was
compared with A.~Rocca's reference volume. When the new box became much smaller,
shell sampling and controlled enlargement of selected parameter ranges were used
to avoid an artificially contracted region. Merged boxes combining both analyses
were also sampled in box-sampling mode to assess the joint admissible set.

\section{Notes on interpretation}
\begin{itemize}
  \item Filtered counts depend strongly on the ratio bounds in
  \texttt{config.py}; relaxing or tightening those bounds changes the accepted
  set dramatically.
  \item Re-sampling densifies local neighbourhoods but can bias coverage; a
  subsequent box-sampling pass on the enclosing box restores a more uniform
  distribution.
  \item Parameter-correlation scatter plots (see
  \texttt{plot\_params\_correlation.py} and the \texttt{multiplot\_*} scripts)
  are useful to visualise possible dependencies in the accepted cloud.
\end{itemize}

\chapter{Stiffness investigation}
\label{chap:stiffness}

\section{Definition used in this work}
A linear constant-coefficient system
\[
\mathbf{y}' = A\mathbf{y} + \mathbf{f}(x)
\]
is considered stiff when all eigenvalues of $A$ have negative real parts and the
stiffness ratio
\[
\frac{|\mathrm{Re}(\overline{\lambda})|}{|\mathrm{Re}(\underline{\lambda})|}
\]
is large, where $\overline{\lambda}$ and $\underline{\lambda}$ are the
eigenvalues with the largest and smallest $|\mathrm{Re}(\lambda)|$,
respectively. Fast transients then coexist with slow ones, which is numerically
demanding for explicit integrators.

Although the iron-homeostasis ODEs are nonlinear, the Jacobian can be evaluated
at a given time point. Locally, the system may then be treated as linear for the
purpose of estimating a stiffness ratio.

\section{Implementation in the code}
In \texttt{model.py}, inside \texttt{diff\_eqs()}, the Jacobian is assembled in
the time window corresponding to the iron perturbation
(approximately $t\in[10000,30000]$\,s, i.e.\ about $2.7$--$8.3$\,h). Eigenvalues
are computed with \texttt{scipy.linalg.eig}. When negative real parts are found
and pairwise stiffness ratios exceed a threshold (e.g.\ $100$), a short report
is written to a stiffness log file.

This diagnostic is optional and intended for model analysis; it does not change
the acceptance criteria of the sampling pipeline.

\section{Practical remarks}
\begin{itemize}
  \item Stiffness, when present, motivates the use of implicit or
  stiff-capable integrators (the default path uses SciPy's \texttt{odeint} /
  LSODA-style switching).
  \item The diagnostic depends on the current parameter set and on the state
  trajectory; a system may look stiff for some accepted parameters and not for
  others.
  \item Extending the model with HIF-related reactions increases the state
  dimension and can alter the eigenvalue spectrum; the same Jacobian-based check
  remains applicable after updating the matrix assembly.
\end{itemize}


\end{document}
